% 14 shorten (shorten, shorten <, shorten >, plain and double arrows) and crossing over: the over-strand has a white halo that interrupts earlier arrows; cube diagram
\documentclass{article}
\pdfcompresslevel=0 \pdfobjcompresslevel=0
\usepackage{amsmath,amssymb}
\usepackage{tikz-cd}
\usetikzlibrary{decorations.pathmorphing}
\pagestyle{empty}
\begin{document}
\begin{tikzcd}[column sep=large]
A \arrow[r,shorten=3pt] & B \arrow[r,shorten <=3pt] & C \arrow[r,shorten >=3pt] & D \arrow[r,shorten <=-2pt,shorten >=-2pt] & E
\end{tikzcd}

\begin{tikzcd}[column sep=large]
A \arrow[r,Rightarrow,shorten=3pt] & B \arrow[r,dash,shorten=1em] & C \arrow[r,hook,two heads,shorten=2mm] & D \arrow[r,equal,shorten >=1ex] & E
\end{tikzcd}

\begin{tikzcd}[column sep=large]
A \arrow[r,shorten=8pt,"f"] \arrow[d,shorten >=4pt,"g"] & B \\ C & D
\end{tikzcd}


\begin{tikzcd}
A \arrow[rd] & B \arrow[ld,crossing over] \\ C & D
\end{tikzcd}

\begin{tikzcd}[row sep=2.5em]
A' \arrow[rr] \arrow[dr,swap,"a"] \arrow[dd,swap,"g"] &&
  B' \arrow[dd,swap,"h" near start] \arrow[dr,"b"] \\
& A \arrow[rr,crossing over,"f" near start] &&
  B \arrow[dd,"k"] \\
C' \arrow[rr,"m" near end] \arrow[dr,swap,"c"] && D' \arrow[dr,swap,"d"] \\
& C \arrow[rr,"l"] \arrow[uu,<-,crossing over,"\beta" near end]&& D
\end{tikzcd}

\begin{tikzcd}
A \arrow[rd] & B \arrow[ld,crossing over,shorten=3pt] \\ C & D \arrow[lu,crossing over,bend left] \arrow[l,crossing over,"f"]
\end{tikzcd}

\begin{tikzcd}[background color=yellow,crossing over clearance=3pt]
A \arrow[rd] \arrow[r,"f" description] & B \arrow[ld,crossing over,Rightarrow] \\ C & D
\end{tikzcd}
\end{document}
